\documentclass[11pt]{article}
\usepackage[margin=1in]{geometry}
\usepackage{knotdiag}

\kdset{scale=0.72, line width=0.9pt, gap=4.2}

\begin{document}

\section*{Conway polynomial of the figure-eight knot}

We use the oriented skein relation
\[
  \skeinrel,
  \qquad
  \nablaknot{unknot}=1,
  \qquad
  \nablaknot{unlink}=0.
\]

Choose a crossing of the figure-eight knot.  Switching it gives an
unknot, while the oriented smoothing gives the negative Hopf link:
\begin{knotcalc}
  \nablaknot{figureeight}
    &= \nablaknot[change=2]{figureeight}
       + z\,\nablaknot{hopf} \\
    &= \nablaknot{unknot}
       + z\,\nablaknot{hopf}.
\end{knotcalc}

Apply the skein relation once more to the Hopf link.  A crossing change
gives the two-component unlink and the smoothing gives an unknot:
\begin{knotcalc}
  \nablaknot{hopf}
    &= \nablaknot{unlinkx}
       - z\,\nablaknot{unknot} \\
    &= \nablaknot{unlink}
       - z\,\nablaknot{unknot} \\
    &= 0-z=-z.
\end{knotcalc}

Therefore
\begin{knotcalc}
  \nablaknot{figureeight}
    &= 1+z(-z) \\
    &= \boxed{1-z^2}.
\end{knotcalc}

\end{document}
